\subsection{Roots of unity}

DEFINITION 2. --- An element \(\zeta\) of \(\mathbf{K}\) is said to be a root of unity if there exists an integer \(n > 0\) such that \(\zeta^n = 1\) ; for every integer \(n > 0\) such that \(\zeta^n = 1\) , \(\zeta\) is called an \(n\) -th root of unity.

It amounts to the same to say that the roots of unity are the elements of finite order of the multiplicative group \(K^{*}\) of non-zero elements of K (I, p.~51). The roots of unity form a subgroup \(\mu_{\infty}(\mathbf{K})\) of \(K^{*}\) , the n-th roots form a subgroup \(\mu_{n}(\mathbf{K})\) of \(\mu_{\infty}(\mathbf{K})\) . We have \(\mu_{\infty}(\mathbf{K}) = \bigcup_{n \geqslant 1} \mu_{n}(\mathbf{K})\) and \(\mu_{n}(\mathbf{K}) \subset \mu_{m}(\mathbf{K})\) if n divides

\begin{enumerate}
\def\labelenumi{\alph{enumi}.}
\setcounter{enumi}{12}
\tightlist
\item
  For every root of unity \(\zeta\) there exists a least integer \(n \geqslant 1\) such that \(\zeta\) belongs to \(\mu_{n}(\mathbf{K})\) , namely the order of \(\zeta\) in the group \(K^{*}\) .
\end{enumerate}

The group \(\mu_n(\mathbf{K})\) , being the set of roots of the polynomial \(\mathbf{X}^n - 1\) , is of finite order \(\leqslant n\) . Let \(p\) be the characteristic of \(\mathbf{K}\) . When \(p = 0\) , or when \(p \neq 0\) and \(n\) is not divisible by \(p\) , the derivative \(n\mathbf{X}^{n-1}\) of \(\mathbf{X}^n - 1\) is relatively prime to \(\mathbf{X}^n - 1\) and so the polynomial \(\mathbf{X}^n - 1\) is then separable; if moreover \(\mathbf{K}\) is algebraically closed, \(\mu_n(\mathbf{K})\) is thus a group of \(n\) elements.

Suppose that K is of non-zero characteristic p and let \(r \geqslant 0\) be an integer; since the mapping \(x \mapsto x^{p^{r}}\) of K into K is injective, we have \(\mu_{np^{r}}(\mathbf{K}) = \mu_n(\mathbf{K})\) for every integer \(n \geqslant 1\) .

We remark that a field may contain no \(n\) -th root of unity other than 1: this is the case for example with the prime fields \(\mathbf{Q}\) and \(\mathbf{F}_2\) for any odd integer \(n\) .

THEOREM 1. --- Let p be the characteristic exponent of K and let n \textgreater{} 0 be an integer. Then the group \(\mu_{n}(\mathbf{K})\) of n-th roots of unity in K is cyclic and its order divides n; when K is algebraically closed and n is prime to p, the group \(\mu_{n}(\mathbf{K})\) is cyclic of order n.

It is enough to prove the first assertion of the theorem which is a consequence of the following more precise lemma:

Lemma 1. --- Let G be a finite subgroup of \(\mathbf{K}^*\) of order \(m\) ; then G is cyclic and we have \(\mathrm{G} = \mu_m(\mathrm{K})\) .

Consider G as Z-module; we have mx = 0 for all \(x \in G\) , hence the annihilator of G is an ideal of the form rZ, where the integer \(r \geqslant 1\) divides m. We thus have \(G \subset \mu_{r}(K)\) . By Lemma 4 (V, p.~74) applied with A = Z and M = G there exists an element x of G of order r; let \(G'\) be the cyclic subgroup of G generated by x. We have \(G' \subset \mu_{r}(K)\) , Card \((G') = r\) and Card \(\mu_{r}(K) \leqslant r\) ; therefore we have

\(G' = \mu_r(K) \supset G\) and \(G\) is cyclic of order \(r\) , hence equal to \(\mu_r(K)\) . Since \(G\) has order \(m\) , we have \(m = r\) and the lemma follows.

PROPOSITION 2. --- Let K be algebraically closed and let p be its characteristic exponent. There exists an isomorphism of \(\mu_{\infty}(\mathbf{K})\) onto the group \(\mathbf{Q}/\mathbf{Z}[1/p]\) .

We have denoted by \(\mathbf{Z}[1/p]\) the subring of \(\mathbf{Q}\) generated by \(1/p\) , that is, the set of rational numbers of the form \(a/p^n\) with \(a \in \mathbf{Z}\) and \(n \geqslant 1\) ; thus we have \(\mathbf{Z}[1/p] = \mathbf{Z}\) if \(\mathbf{K}\) is of characteristic 0.

Let \((\nu_{n})_{n\geqslant 1}\) be the strictly increasing sequence consisting of all the integers which are not divisible by \(p\) , if \(p\neq 1\) ; put \(\lambda_{n} = \nu_{1}\nu_{2}\dots \nu_{n}\) and denote by \(\mathrm{H}_n\) the group of \(\lambda_{n}\) -th roots of unity; we have \(\mathrm{H}_{n + 1}\supset \mathrm{H}_n\) and \(\mu_{\infty}(\mathbf{K}) = \cup_{n}\mathrm{H}_{n}\) . Since \(\mathrm{H}_n\) is cyclic of order \(\lambda_{n}\) (Th. 1), there exists a sequence

\((\alpha_{n})_{n\geqslant 1}\) of roots of unity such that \(\alpha_{n}\) generates \(\mathrm{H}_n\) and \(\alpha_{n} = \alpha_{n + 1}^{\nu_{n + 1}}\)

Further let \(\beta_{n}\) be the residue class mod \(\mathbf{Z}[1/p]\) of \(1/\lambda_{n}\) and let \(\mathbf{H}_{n}^{\prime}\) be the cyclic subgroup of \(\mathbf{Q}/\mathbf{Z}[1/p]\) generated by \(\beta_{n}\) . It is clear that \(\beta_{n} = \nu_{n+1}\beta_{n+1}\) and \(\mathbf{H}_{n}^{\prime}\) is of order \(\lambda_{n}\) because \(\lambda_{n}\) is not divisible by \(p\) if \(p \neq 1\) . Thus there exists for all \(n \geqslant 1\) an isomorphism \(\varphi_{n}: \mathrm{H}_{n} \to \mathrm{H}_{n}^{\prime}\) such that \(\varphi_{n}(\alpha_{n}) = \beta_{n}\) and the relations \(\alpha_{n} = \alpha_{n+1}^{\nu_{n+1}}\) , \(\beta_{n} = \nu_{n+1}\beta_{n+1}\) show that \(\varphi_{n+1}\) extends \(\varphi_{n}\) . Finally there exists a unique isomorphism \(\varphi\) of \(\mu_{\infty}(\mathbf{K})\) onto \(\mathbf{Q}/\mathbf{Z}[1/p]\) extending the isomorphisms \(\varphi_{n}\) , that is, \(\varphi(\alpha_{n}) = \beta_{n}\) for all \(n \geqslant 1\) .

Remarks. --- 1) When K is an algebraically closed field of characteristic 0, the group \(\mu_{\infty}(\mathbf{K})\) is thus isomorphic (not canonically) to \(\mathbf{Q}/\mathbf{Z}\) . * When K is the field \(\mathbf{C}\) of complex numbers, we can write down an explicit such isomorphism; in effect the mapping \(x \mapsto e^{2\pi ix}\) is a homomorphism of \(\mathbf{Q}\) into \(\mathbf{C}^*\) with kernel \(\mathbf{Z}\) and image \(\mu_{\infty}(\mathbf{C})\) ; thus by passage to quotients it defines an isomorphism of \(\mathbf{Q}/\mathbf{Z}\) onto \(\mu_{\infty}(\mathbf{C})\) . *

\begin{enumerate}
\def\labelenumi{\arabic{enumi})}
\setcounter{enumi}{1}
\item
  The following result may be shown (cf.~V, p.~165, Ex. 21): let G and H be two commutative groups all of whose elements are of finite order. Suppose that for every integer \(n \geqslant 1\) , the equation \(nx = 0\) has the same number of solutions, assumed finite in G and in H. Then the groups G and H are isomorphic. This provides a new proof of Prop. 2.
\item
  For every prime number \(l\) put \(\mu_{l^{\infty}}(\mathbf{K}) = \bigcup_{n\geqslant 0}\mu_{l^{n}}(\mathbf{K})\) . When \(l\) is the characteristic
\end{enumerate}

\(p\) of \(\mathbf{K}\) , we have \(\mu_{p^{\infty}}(\mathbf{K}) = \{1\}\) . From I, p.~80, Th. 4, we deduce that \(\mu_{\infty}(\mathbf{K})\) is a direct sum of the subgroups \(\mu_{l^{\infty}}(\mathbf{K})\) , where \(l\) runs over the set of all prime numbers distinct from \(p\) . For a given prime number \(l\) only two cases are possible: either \(\mu_{l^{\infty}}(\mathbf{K})\) is finite and then \(\mu_{l^{\infty}}(\mathbf{K})\) is isomorphic to \(\mathbf{Z}/l^{n}\mathbf{Z}\) for a suitable \(n\) (Th. 1), or \(\mu_{l^{\infty}}(\mathbf{K})\) is infinite and then \(\mu_{l^{\infty}}(\mathbf{K})\) is isomorphic to \(\mathbf{Z}[l^{-1}]/\mathbf{Z}\) (cf.~Remark 2).

